\documentclass[12pt, a4paper]{article}

% -------------------------------------------------------
% Encodage et langue
% -------------------------------------------------------
\usepackage[utf8]{inputenc}
\usepackage[T1]{fontenc}
\usepackage[french]{babel}

% -------------------------------------------------------
% Mise en page
% -------------------------------------------------------
\usepackage[top=2.5cm, bottom=2.5cm, left=2.5cm, right=2.5cm]{geometry}
\usepackage{parskip}          % espace entre paragraphes, pas d'indentation
\usepackage{microtype}        % meilleure typographie

% -------------------------------------------------------
% Mathématiques
% -------------------------------------------------------
\usepackage{amsmath}
\usepackage{amssymb}

% -------------------------------------------------------
% Tableaux
% -------------------------------------------------------
\usepackage{booktabs}
\usepackage{array}
\usepackage{tabularx}

% -------------------------------------------------------
% Couleurs et boîtes
% -------------------------------------------------------
\usepackage{xcolor}
\usepackage{tcolorbox}
\tcbuselibrary{skins, breakable}

% Couleurs personnalisées
\definecolor{bleuCours}{RGB}{30, 90, 160}
\definecolor{fondNote}{RGB}{240, 248, 255}
\definecolor{fondHumour}{RGB}{255, 249, 230}
\definecolor{bordHumour}{RGB}{255, 180, 0}
\definecolor{fondAttention}{RGB}{255, 243, 243}
\definecolor{bordAttention}{RGB}{200, 50, 50}

% Boîte pour les notes humoristiques
\newtcolorbox{humourbox}{
  colback=fondHumour,
  colframe=bordHumour,
  fontupper=\itshape,
  left=6pt, right=6pt, top=4pt, bottom=4pt,
  arc=4pt
}

% Boîte pour les avertissements
\newtcolorbox{attentionbox}{
  colback=fondAttention,
  colframe=bordAttention,
  fontupper=\itshape,
  left=6pt, right=6pt, top=4pt, bottom=4pt,
  arc=4pt
}

% Boîte pour les résumés / étapes
\newtcolorbox{resumebox}{
  colback=fondNote,
  colframe=bleuCours,
  left=6pt, right=6pt, top=6pt, bottom=6pt,
  arc=4pt
}

% -------------------------------------------------------
% Titres
% -------------------------------------------------------
\usepackage{titlesec}
\titleformat{\section}{\large\bfseries\color{bleuCours}}{}{0em}{}[\titlerule]
\titleformat{\subsection}{\normalsize\bfseries\color{bleuCours}}{}{0em}{}

% -------------------------------------------------------
% En-tête et pied de page
% -------------------------------------------------------
\usepackage{fancyhdr}
\pagestyle{fancy}
\fancyhf{}
\rhead{\textcolor{bleuCours}{\textit{Statistiques — Test du Khi carré}}}
\lhead{\textcolor{bleuCours}{\textit{Cours}}}
\rfoot{\thepage}
\renewcommand{\headrulewidth}{0.4pt}

% -------------------------------------------------------
% Commande raccourci
% -------------------------------------------------------
\newcommand{\khi}{\chi^2}

% -------------------------------------------------------
% Document
% -------------------------------------------------------
\begin{document}

% Page de titre
\begin{center}
  {\LARGE \textbf{\textcolor{bleuCours}{Le Test du Khi carré} ($\chi^2$)}}\\[0.8em]
  {\large Cours de statistiques}\\[0.4em]
  {\normalsize \today}
\end{center}

\vspace{0.5em}

\begin{humourbox}
« Le Khi carré, c'est comme votre beau-frère qui vérifie si vos dés sont pipés ---
sauf que lui, il a une formule. »
\end{humourbox}

\vspace{1em}

% ===================================================
\section{1. Introduction}
% ===================================================

Le \textbf{test du Khi carré} (noté $\chi^2$) est un test statistique utilisé pour
déterminer s'il existe une association significative entre deux variables catégorielles,
ou si une distribution observée s'écarte significativement d'une distribution théorique
attendue.

Autrement dit : est-ce que ce qu'on \textbf{observe} dans la vraie vie ressemble à ce
qu'on \textbf{s'attendait} à voir ? Si la réponse est « pas du tout », les données ont
quelque chose à nous dire.

Il existe deux grandes variantes :

\vspace{0.5em}
\begin{center}
\begin{tabular}{lp{8cm}}
  \toprule
  \textbf{Variante} & \textbf{Usage} \\
  \midrule
  Test d'ajustement (\textit{goodness of fit}) & Compare une distribution observée à
    une distribution théorique \\[4pt]
  Test d'indépendance & Vérifie si deux variables catégorielles sont indépendantes \\
  \bottomrule
\end{tabular}
\end{center}

% ===================================================
\section{2. Conditions d'application}
% ===================================================

\begin{attentionbox}
Avant de sortir votre $\chi^2$, assurez-vous que vos données sont invitées au bon party.
\end{attentionbox}

\vspace{0.5em}

Avant d'appliquer le test du Khi carré, il faut vérifier les conditions suivantes :

\begin{itemize}
  \item Les données sont des \textbf{fréquences} (effectifs), non des pourcentages ou
        des moyennes.
  \item Les observations sont \textbf{indépendantes} les unes des autres.
  \item Chaque effectif \textbf{théorique} doit être $\geq 5$ dans au moins 80\,\% des
        cellules.
  \item Aucun effectif théorique ne doit être \textbf{inférieur à 1}.
\end{itemize}

% ===================================================
\section{3. La statistique du Khi carré}
% ===================================================

\subsection{Formule générale}

\begin{equation}
  \chi^2 = \sum_{i} \frac{(O_i - E_i)^2}{E_i}
\end{equation}

Où :
\begin{itemize}
  \item $O_i$ = effectif \textbf{observé} dans la catégorie $i$
  \item $E_i$ = effectif \textbf{attendu} (théorique) dans la catégorie $i$
  \item $\sum$ = somme sur toutes les cellules
\end{itemize}

\begin{humourbox}
Traduction : on regarde combien la réalité s'éloigne de la théorie, on met ça au carré
pour que personne ne se cache derrière un signe négatif, et on additionne le tout.
\end{humourbox}

% ===================================================
\section{4. Test d'ajustement (\textit{Goodness of Fit})}
% ===================================================

\subsection{Objectif}

Vérifier si une distribution observée correspond à une distribution théorique connue
(ex. : distribution uniforme, distribution normale, etc.).

\subsection{Exemple}

\begin{humourbox}
Imaginez que vous soupçonnez votre ami de tricher aux jeux de société. Le Khi carré va
régler ça scientifiquement --- et diplomatiquement.
\end{humourbox}

\vspace{0.5em}

Un dé à 6 faces est lancé 120 fois. On observe :

\vspace{0.5em}
\begin{center}
\begin{tabular}{lcccccc}
  \toprule
  \textbf{Face} & 1 & 2 & 3 & 4 & 5 & 6 \\
  \midrule
  \textbf{Observé} ($O_i$) & 18 & 22 & 17 & 25 & 19 & 19 \\
  \textbf{Attendu}  ($E_i$) & 20 & 20 & 20 & 20 & 20 & 20 \\
  \bottomrule
\end{tabular}
\end{center}

\vspace{0.5em}

\textbf{Calcul :}

\begin{align*}
  \chi^2 &= \frac{(18-20)^2}{20} + \frac{(22-20)^2}{20} + \frac{(17-20)^2}{20}
            + \frac{(25-20)^2}{20} + \frac{(19-20)^2}{20} + \frac{(19-20)^2}{20} \\
         &= \frac{4}{20} + \frac{4}{20} + \frac{9}{20} + \frac{25}{20}
            + \frac{1}{20} + \frac{1}{20}
          = \frac{44}{20} = 2{,}2
\end{align*}

\textbf{Degrés de liberté :} $dl = k - 1 = 6 - 1 = 5$

\begin{humourbox}
Résultat : $2{,}2 < 11{,}070$. Le dé est innocent. Votre ami peut dormir tranquille.
\end{humourbox}

% ===================================================
\section{5. Test d'indépendance}
% ===================================================

\subsection{Objectif}

Déterminer si deux variables catégorielles sont \textbf{indépendantes} ou
\textbf{associées}.

\subsection{Tableau de contingence}

Les données sont organisées dans un \textbf{tableau de contingence} à $r$ lignes et
$c$ colonnes.

\textbf{Effectif attendu pour chaque cellule :}

\begin{equation}
  E_{ij} = \frac{(\text{Total ligne } i) \times (\text{Total colonne } j)}
               {\text{Total général}}
\end{equation}

\textbf{Degrés de liberté :}

\begin{equation}
  dl = (r - 1)(c - 1)
\end{equation}

\subsection{Exemple}

On veut savoir si le sexe d'un individu est lié à sa préférence pour un sport.

\vspace{0.5em}
\begin{center}
\begin{tabular}{lcccc}
  \toprule
   & \textbf{Football} & \textbf{Basketball} & \textbf{Tennis} & \textbf{Total} \\
  \midrule
  \textbf{Hommes} & 30 & 15 & 5  & \textbf{50}  \\
  \textbf{Femmes} & 10 & 25 & 15 & \textbf{50}  \\
  \textbf{Total}  & \textbf{40} & \textbf{40} & \textbf{20} & \textbf{100} \\
  \bottomrule
\end{tabular}
\end{center}

\vspace{0.5em}

\textbf{Calcul des effectifs attendus :}

\[
  E_{11} = \frac{50 \times 40}{100} = 20 \quad
  E_{12} = \frac{50 \times 40}{100} = 20 \quad
  E_{13} = \frac{50 \times 20}{100} = 10
\]
\[
  E_{21} = 20 \quad E_{22} = 20 \quad E_{23} = 10
\]

\textbf{Calcul du $\chi^2$ :}

\begin{align*}
  \chi^2 &= \frac{(30-20)^2}{20} + \frac{(15-20)^2}{20} + \frac{(5-10)^2}{10}
            + \frac{(10-20)^2}{20} + \frac{(25-20)^2}{20} + \frac{(15-10)^2}{10} \\
         &= 5 + 1{,}25 + 2{,}5 + 5 + 1{,}25 + 2{,}5 = 17{,}5
\end{align*}

\textbf{Degrés de liberté :} $dl = (2-1)(3-1) = 2$

% ===================================================
\section{6. Prise de décision}
% ===================================================

\begin{humourbox}
Le statisticien joue les juges : $H_0$ est présumée innocente jusqu'à preuve du
contraire. Votre $\chi^2$ est l'avocat de l'accusation.
\end{humourbox}

\subsection{Hypothèses}

\begin{itemize}
  \item $H_0$ (hypothèse nulle) : Il n'y a \textbf{pas} d'association entre les variables
        (ou la distribution observée suit la distribution théorique).
  \item $H_1$ (hypothèse alternative) : Il \textbf{existe} une association entre les
        variables.
\end{itemize}

\subsection{Règle de décision}

On compare la valeur calculée de $\chi^2$ à la \textbf{valeur critique}
$\chi^2_{\alpha,\,dl}$ issue de la table du Khi carré pour un seuil de signification
$\alpha$ (généralement 0,05).

\vspace{0.5em}
\begin{center}
\begin{tabular}{lp{7cm}}
  \toprule
  \textbf{Condition} & \textbf{Décision} \\
  \midrule
  $\chi^2_{\text{calc}} > \chi^2_{\text{critique}}$ &
    \textbf{Rejeter $H_0$} $\rightarrow$ association significative \\[4pt]
  $\chi^2_{\text{calc}} \leq \chi^2_{\text{critique}}$ &
    \textbf{Ne pas rejeter $H_0$} $\rightarrow$ pas d'association significative \\
  \bottomrule
\end{tabular}
\end{center}

\vspace{0.3em}
On peut aussi comparer la \textbf{p-valeur} au seuil $\alpha$ : si p-valeur $< \alpha$,
on rejette $H_0$.

\subsection{Table des valeurs critiques (extrait)}

\vspace{0.5em}
\begin{center}
\begin{tabular}{cccc}
  \toprule
  \textbf{dl} & $\alpha = 0{,}10$ & $\alpha = 0{,}05$ & $\alpha = 0{,}01$ \\
  \midrule
  1 & 2,706 & 3,841 &  6,635 \\
  2 & 4,605 & 5,991 &  9,210 \\
  3 & 6,251 & 7,815 & 11,345 \\
  4 & 7,779 & 9,488 & 13,277 \\
  5 & 9,236 & 11,070 & 15,086 \\
  \bottomrule
\end{tabular}
\end{center}

\subsection{Reprise des exemples}

\textbf{Exemple 1 (dé) :} $\chi^2 = 2{,}2$ ; $dl = 5$ ;
$\chi^2_{\text{critique}}$ ($\alpha = 0{,}05$) $= 11{,}070$

$\Rightarrow\ 2{,}2 < 11{,}070\ \Rightarrow$ \textbf{On ne rejette pas $H_0$} :
le dé semble équitable.

\vspace{0.5em}

\textbf{Exemple 2 (sport/sexe) :} $\chi^2 = 17{,}5$ ; $dl = 2$ ;
$\chi^2_{\text{critique}}$ ($\alpha = 0{,}05$) $= 5{,}991$

$\Rightarrow\ 17{,}5 > 5{,}991\ \Rightarrow$ \textbf{On rejette $H_0$} :
il y a une association significative entre le sexe et la préférence sportive.

% ===================================================
\section{7. Résumé --- Démarche en 5 étapes}
% ===================================================

\begin{humourbox}
Cinq étapes. Pas six. Pas quatre. Cinq. Vous pouvez le faire.
\end{humourbox}

\vspace{0.5em}

\begin{resumebox}
\begin{enumerate}
  \item Formuler $H_0$ et $H_1$
  \item Choisir le seuil de signification $\alpha$ (ex. : 0,05)
  \item Calculer $\chi^2$ et les degrés de liberté
  \item Trouver la valeur critique dans la table
  \item Prendre la décision et interpréter
\end{enumerate}
\end{resumebox}

\vspace{0.5em}

\begin{humourbox}
Félicitations, vous êtes maintenant capables de débusquer les associations cachées dans
n'importe quel tableau. Utilisez ce pouvoir avec sagesse (et avec $\alpha = 0{,}05$).
\end{humourbox}

% ===================================================
\section{8. Limites du test}
% ===================================================

\begin{humourbox}
Même les super-héros ont des faiblesses. Le Khi carré aussi.
\end{humourbox}

\vspace{0.5em}

\begin{itemize}
  \item Ne s'applique \textbf{pas} aux données continues (sauf si regroupées en classes).
  \item Ne mesure \textbf{pas la force} de l'association (pour cela, utiliser le
        \textbf{V de Cramér}).
  \item Sensible à la \textbf{taille de l'échantillon} : un très grand échantillon peut
        rejeter $H_0$ même pour une association triviale.
\end{itemize}

\begin{humourbox}
En d'autres mots : avec 10\,000 observations, même « les gens qui mangent des céréales
au petit-déjeuner sont légèrement plus grands de 0,2~cm » devient statistiquement
significatif. Ce n'est pas la même chose qu'être \textbf{intéressant}.
\end{humourbox}

\subsection{V de Cramér (mesure d'effet)}

\begin{equation}
  V = \sqrt{\frac{\chi^2}{n \times \min(r-1,\; c-1)}}
\end{equation}

\begin{itemize}
  \item $V$ proche de 0 $\rightarrow$ association faible
  \item $V$ proche de 1 $\rightarrow$ association forte
\end{itemize}

\end{document}
